% Reader-first, claim-preserving front matter for WOW-284.
% This fragment is not included automatically.  It rewrites the prior v1-paper
% opening without promoting any post-v1 theorem whose audit gate remains open.

\begin{abstract}
WOW-284 predicts that the minimum dual degree of every connected graph of order
at least three and girth at least five is at most the negative of its least
distance eigenvalue.  The degree-seven Hoffman--Singleton graph disproves the
conjecture.  More generally, a degree-$k$ Moore graph of diameter two satisfies
\[
  \delta^*(G)=k,
  \qquad
  \lambda_{\min}(D(G))=-\frac{3+\sqrt{4k-3}}2,
\]
so WOW-284 holds exactly for $k\le3$, with equality at $k=3$, and fails for every
realisable $k>3$.  We give a self-contained coordinate certificate for the
Hoffman--Singleton example and exact induced counterexamples of orders
$38,39,40$, and $42$.  For every connected $k$-regular graph of girth at least
five and diameter three we also prove
\[
  D=3J+(k-3)I-2A-A^2,
\]
which converts WOW-284 into an exact shifted adjacency-spectrum condition.  All
finite spectral decisions use exact arithmetic.  Lean~4.31 kernel-checks the
complete $50$-vertex proof and finite spectral certificates for the four smaller
constructions.
\end{abstract}

\section*{Introduction}

Let $G$ be a finite simple connected graph.  Write $D(G)$ for its distance
matrix and $\lambda_{\min}(D(G))$ for the least distance eigenvalue.  For a
vertex $v$, define
\[
  d^*(v)=\frac1{d(v)}\sum_{u\in N(v)}d(u),
  \qquad
  \delta^*(G)=\min_{v\in V(G)}d^*(v).
\]
Aouchiche and Hansen record the following Graffiti conjecture, attributed to
Fajtlowicz's \emph{Written on the Wall} report.

\begin{conjecture}[WOW-284]\label{conj:wow284-reader-first}
If $G$ is connected, has at least three vertices, and has girth at least five,
then
\[
  \delta^*(G)\le-\lambda_{\min}(D(G)).
\]
\end{conjecture}

The conjecture is false for a simple reason.  In a degree-$k$ Moore graph of
diameter two, adjacent vertices have no common neighbour and nonadjacent
vertices have exactly one.  If $A$, $I$, and $J$ denote the adjacency, identity,
and all-ones matrices, then
\[
  A^2=(k-1)I-A+J,
  \qquad
  D=2J-2I-A.
\]
On the orthogonal complement of the all-ones vector, the nonprincipal adjacency
eigenvalues are the roots of
\[
  x^2+x-(k-1)=0.
\]
It follows that
\[
  \lambda_{\min}(D)=-\frac{3+\sqrt{4k-3}}2.
\]
Regularity gives $\delta^*(G)=k$, and therefore the conjectured inequality is
equivalent to
\[
  k\le\frac{3+\sqrt{4k-3}}2.
\]
For integers $k\ge2$, this holds exactly when $k\le3$.  The Petersen graph is
the equality case $k=3$, whereas the degree-seven Hoffman--Singleton graph has
\[
  \delta^*=7,
  \qquad
  \lambda_{\min}(D)=-4,
\]
and hence violates WOW-284 by the strict gap $3$.

This paper develops that observation in three directions.  First, we give a
fully labelled coordinate construction of the Hoffman--Singleton graph and
verify the common-neighbour identity directly.  This makes the disproof
self-contained and supplies exact graph data for independent checking.  Second,
we extract smaller induced counterexamples of orders $38,39,40$, and $42$.
The $40$- and $42$-vertex examples are regular; the $38$- and $39$-vertex
examples show that the phenomenon persists after small asymmetric deletions.
Third, we isolate the operator mechanism behind the regular diameter-three
examples.  If $G$ is connected, $k$-regular, has girth at least five, and has
diameter three, then
\[
  D=3J+(k-3)I-2A-A^2.
\]
Thus every nonprincipal adjacency eigenvalue $\theta$ produces the distance
eigenvalue
\[
  k-2-(\theta+1)^2,
\]
and $G$ is a strict counterexample precisely when all nonprincipal adjacency
eigenvalues lie in the open interval centred at $-1$ with radius
$\sqrt{2k-2}$.

The distance spectra of the classical Moore graphs and the general
minimal-cage distance-polynomial framework are established in the literature.
Our contribution is to make the connection with WOW-284 explicit, record the
sharp Moore-degree threshold, provide self-contained exact certificates, and
organise the smaller constructions through a common spectral mechanism.  We do
not claim that the displayed classical distance spectra are new, that the
$38$-vertex example is globally minimal, or that a targeted literature search
settles priority over unpublished observations.

The proof is organised accordingly.  Section~\ref{sec:moore} gives the short
Moore-graph calculation.  The next section supplies the coordinate certificate.
We then derive the smaller induced examples, prove the general diameter-three
identity, and treat Moore second subconstituents.  Exact characteristic
polynomials, Sturm certificates, rational $LDL^{\mathsf T}$ decompositions, and
formal-verification details are recorded only where they certify a stated
mathematical step.

\phantomsection\label{thm:moore-wow-reader-first}
\begin{theoremA}
Let $G$ be a Moore graph of diameter two and degree $k\ge2$.  Then
\[
  |V(G)|=k^2+1,
  \qquad
  g(G)=5,
  \qquad
  \delta^*(G)=k,
  \qquad
  \lambda_{\min}(D(G))=-\frac{3+\sqrt{4k-3}}2.
\]
Consequently, $G$ satisfies WOW-284 if and only if $k\le3$, with equality
exactly when $k=3$.
\end{theoremA}
